\documentclass[12pt]{article}
\usepackage[utf8]{inputenc} % Кодировка UTF-8 (pdflatex)
\usepackage[T2A]{fontenc}    % Русские шрифты
\usepackage[russian]{babel}  % Русский язык для babel
\usepackage{graphicx}        % Графика
\usepackage{microtype}       % Микротипографика
\usepackage{amssymb}         % Математические символы
\usepackage{amsmath}         % Математика
\usepackage{tikz}            % Рисование
\usepackage{circuitikz}      % Электрические схемы
\usepackage{fix-cm}          % Чтобы работал большой размер шрифта
\usepackage{anyfontsize}     % Для любого размера шрифта

% Заданные раздельные отступы страницы (для pdflatex)
% левый 30мм, правый 20мм, верхний 25мм, нижний 25мм
\usepackage[a4paper,left=20mm,right=20mm,top=20mm,bottom=20mm]{geometry}

% Для pdflatex: используем стандартные пакеты шрифтов
% Если нужно Times-подобие: раскомментируйте одну из следующих строк
% \usepackage{mathptmx}
% \usepackage{newtxtext,newtxmath}

\title{ТА - СРС8 - ЭН-25 Андреев В1}
\author{}
\date{}

\begin{document}

\maketitle

\section*{Простое Задание}

Дано две группы и условие. 
Если взять $\forall$ пару студентов внутри одной группы, 
то гарантировано, что хотя бы один из них делал работу 
не сам. Соответственно, условие всегда будет выполняться,
только если в группе только один студент сделал работу сам.
Таким образом, студентов, которые сделали работу сами, не больше двух, 
исходя из количества групп.

\section*{Сложное Задание}

\begin{center}
% Увеличим расстояние между колонками таблицы (пространство между полями)
\setlength{\tabcolsep}{8mm}
\begin{tabular}{c c}
\begin{tikzpicture}[scale=0.8,baseline=(current bounding box.center)]
	% размеры поля
	\def\cols{8}
	\def\rows{8}

	% Стандартный цвет клетки (если цвет для конкретной ячейки не задан)
	\def\defaultcell{cyan!50}

	% Макрос для задания цвета конкретной ячейки первого поля: \cellA{x}{y}{<цвет>}
	\newcommand{\cellA}[3]{\expandafter\gdef\csname cellcolorA@#1@#2\endcsname{#3}}

	% Пример заданных ячеек первого поля (вы можете добавить свои \cellA{<x>}{<y>}{<color>} перед рисунком)
	\cellA{0}{0}{blue!70}

	\cellA{3}{0}{blue!70}
	\cellA{5}{1}{blue!70}
	\cellA{5}{2}{blue!70}
	\cellA{6}{3}{blue!70}
	\cellA{0}{2}{blue!70}
	\cellA{1}{3}{blue!70}
	\cellA{7}{4}{blue!70}
	\cellA{7}{5}{blue!70}
	\cellA{6}{5}{blue!70}
	\cellA{6}{4}{blue!70}
	\cellA{5}{5}{blue!70}
	\cellA{4}{5}{blue!70}
	\cellA{3}{5}{blue!70}
	\cellA{4}{4}{blue!70}

	\cellA{0}{7}{blue!70}
	\cellA{0}{6}{blue!70}
	\cellA{1}{7}{blue!70}
	\cellA{3}{7}{blue!70}

	\cellA{7}{7}{blue!70}

	% Рисуем клетки: используем заданный цвет, если он есть, иначе - стандартный
	% Заполняем чуть меньше каждой клетки и с прозрачностью, чтобы промежутки были белыми
	\foreach \x in {0,...,7}{
		\foreach \y in {0,...,7}{
			\ifcsname cellcolorA@\x@\y\endcsname
				\edef\c{\csname cellcolorA@\x@\y\endcsname}
				\fill[\c,fill opacity=0.5] (\x+0.05,\y+0.05) rectangle ++(0.9,0.9);
			\else
				\fill[\defaultcell,fill opacity=0.12] (\x+0.05,\y+0.05) rectangle ++(0.9,0.9);
			\fi
		}
	}

	% рамка вокруг поля
	\draw[line width=1pt] (0,0) rectangle (\cols,\rows);
\end{tikzpicture} & \begin{tikzpicture}[scale=0.8,baseline=(current bounding box.center)]
	% размеры поля
	\def\cols{8}
	\def\rows{8}

	% Стандартный цвет клетки (если цвет для конкретной ячейки не задан)
	\def\defaultcell{cyan!50}

	% Макрос для задания цвета конкретной ячейки второго поля: \cellB{x}{y}{<цвет>}
	\newcommand{\cellB}[3]{\expandafter\gdef\csname cellcolorB@#1@#2\endcsname{#3}}

	% Пример заданных ячеек второго поля (добавляйте \cellB{<x>}{<y>}{<color>} перед рисунком)
	% Закрашенные клетки

	% y = 0
	\cellB{0}{0}{blue!70}
	\cellB{1}{0}{blue!70}
	\cellB{2}{0}{blue!70}
	\cellB{4}{0}{blue!70}
	\cellB{7}{0}{blue!70}
	% y = 1
	\cellB{4}{1}{blue!70}

	% y = 2
	\cellB{5}{2}{blue!70}
	\cellB{6}{2}{blue!70}

	% y = 3
	\cellB{0}{3}{blue!70}
	\cellB{6}{3}{blue!70}
	% y = 4
	\cellB{0}{4}{blue!70}
	\cellB{3}{4}{blue!70}
	\cellB{4}{4}{blue!70}

	% y = 5
	\cellB{0}{5}{blue!70}
	\cellB{3}{5}{blue!70}
	\cellB{4}{5}{blue!70}

	% y = 6
	\cellB{3}{6}{blue!70}
	\cellB{5}{6}{blue!70}
	\cellB{1}{6}{blue!70}
	\cellB{2}{6}{blue!70}

	% y = 7
	\cellB{1}{7}{blue!70}
	\cellB{2}{7}{blue!70}
	\cellB{7}{7}{blue!70}

	% Рисуем клетки: используем заданный цвет, если он есть, иначе - стандартный
	% Заполняем чуть меньше каждой клетки и с прозрачностью, чтобы промежутки были белыми
	\foreach \x in {0,...,7}{
		\foreach \y in {0,...,7}{
			\ifcsname cellcolorB@\x@\y\endcsname
				\edef\c{\csname cellcolorB@\x@\y\endcsname}
				\fill[\c,fill opacity=0.5] (\x+0.05,\y+0.05) rectangle ++(0.9,0.9);
			\else
				\fill[\defaultcell,fill opacity=0.12] (\x+0.05,\y+0.05) rectangle ++(0.9,0.9);
			\fi
		}
	}

	% рамка вокруг поля
	\draw[line width=1pt] (0,0) rectangle (\cols,\rows);
\end{tikzpicture} \\
Изначальное поле & Следующая итерация
\end{tabular}
\end{center}


\end{document}